\section{Related Work}
\label{sec:related}

\subsection{Event-based infrared hardware}
The DARPA FENCE programme funds event-based infrared focal plane arrays with
embedded neuromorphic processing below 1.5\,W~\cite{fence}. Raytheon
demonstrated a first-of-its-kind event-based MWIR camera in April 2026 in
ground-vehicle, aircraft, and live-fire scenes~\cite{raytheon2026}. On the
device side, uncooled Poisson-bolometer LWIR pixels based on spintronic
stochastic switching achieve nanosecond response at sub-\si{\micro\watt}
power~\cite{mousa2026poisson}, and neural-bolometer architectures embed
in-pixel neuromorphic compute~\cite{neuralbolometers2026}. None of these
efforts has released data.

\subsection{Novelty audit}
We systematically checked for existing thermal event data as of July 2026.
EVIF~\cite{geng2024evif} pairs \emph{visible} DAVIS346 events with thermal
frames and synthesizes events only from the visible channel.
DVS-PedX~\cite{sakhai2026dvspedx} converts JAAD \emph{RGB} driving video
with v2e---the methodological precedent for our work, but explicitly
visible-light. SEVD~\cite{sevd} renders synthetic visible multi-view events
in CARLA. The CVPR EventVision 2025 workshop accepted zero
thermal/infrared papers, and the PBVS thermal workshop has no event track.
A patent on event-based thermal camera systems exists~\cite{wo2022patent}
but does not affect publication. To our knowledge, ours is the first public
synthetic thermal event stream data and benchmark.

\subsection{Event simulators}
ESIM~\cite{rebecq2018esim} introduced adaptive rendering; vid2e/v2e
established recycling annotated video with realistic noise
models~\cite{gehrig2020vid2e,hu2021v2e}; DVS-Voltmeter models Brownian
threshold processes~\cite{lin2022voltmeter}; ADV2E decouples analog pixel
filtering from source frame rate~\cite{jiang2024adv2e}; PECS adds physical
optics and multispectral rendering~\cite{han2024pecs}; V2CE learns the
mapping~\cite{zhang2024v2ce}; V2V outputs voxels directly~\cite{lou2025v2v}.
All target visible-light DVS statistics; none models microbolometer
physics or AGC.

\subsection{Event-based detection}
RED~\cite{perot2020red} and RVT~\cite{gehrig2023rvt} set the recurrent
standard; SAST~\cite{peng2024sast}, state-space models~\cite{zubic2024ssm},
and label-efficient LEOD~\cite{wu2024leod} extend it. Thermal-frame
detection is comparatively mature (e.g.\ FLIR
benchmarks~\cite{fliradas}). Cross-modal sim-to-real gaps are documented
by~\cite{stoffregen2020sim2real}, and event fidelity can be scored by
EQS~\cite{chanda2025eqs}. We build on these baselines and metrics while
isolating the \emph{input modality} as the experimental variable.

\subsection{Thermal datasets}
KAIST~\cite{hwang2015kaist} provides aligned visible--thermal street video;
FLIR ADAS v2~\cite{fliradas} provides annotated thermal driving frames and,
critically for us, a continuous video test set; HIT-UAV~\cite{suo2023hituav}
provides UAV thermal imagery under a permissive licence;
LLVIP~\cite{jia2021llvip} and BIRDSAI~\cite{bondi2020birdsai} cover
surveillance and aerial thermal video. None provides events.
